
This section describes the hardware, software, workload and timing
procedure behind every measurement in the paper, and states what our
records do not contain. The full provenance audit, with a
file-level source for each statement below, is part of the released
repository (see Code and Data Availability).

\subsection{Hardware}
\label{sec:meth-hardware}

We compare three backends running the same AlphaFold2 inference code
(Table~\ref{tab:meth-hardware}). CPU and GPU runs used Google Colab
runtimes; TPU runs used a Google Kubernetes Engine cluster in
\texttt{us-west4} provided through Stanford's ME344 course.

\begin{table}[H]
\centering\small
\caption{Hardware and software per backend. ``Not recorded'' means our
run artifacts do not contain the value; these gaps are discussed in
Section~\ref{sec:limitations}.}
\label{tab:meth-hardware}
\begin{tabularx}{\linewidth}{>{\raggedright\arraybackslash}p{0.17\linewidth}>{\raggedright\arraybackslash}X>{\raggedright\arraybackslash}X>{\raggedright\arraybackslash}X}
\toprule
 & CPU (Colab) & GPU (Colab) & TPU (GKE) \\
\midrule
Device & x86\_64 CPU runtime, 2 vCPU (from our write-up, not a
recorded field in the run artifacts); August host CPU model not
recorded; September reruns: Intel Xeon @ 2.20\,GHz (family 6,
model 79) & NVIDIA Tesla T4, 15{,}360\,MiB, compute capability 7.5,
driver 580.82.07 & TPU v5e slice (\texttt{tpu-v5-lite-podslice},
topology 2$\times$4): 8 chips, 16.9\,GB HBM limit per chip \\
Host CPU / RAM & same as device & not recorded (August) & not recorded \\
JAX & unpinned in August (resolved version not recorded); 0.10.2 in
the September reruns & unpinned, \texttt{jax[cuda12]} (resolved
version not recorded) & \texttt{jax[tpu]==0.10.2} \\
Scheduling & interactive notebook & interactive notebook & Kubernetes
Job, Kueue queue, 8 chips requested \\
\bottomrule
\end{tabularx}
\end{table}

Two properties of this setup bear on the results. The TPU allocation
is always the full 8-chip slice, but in the default single-query
execution path only one chip performs work: seven of the eight chips
report zero bytes of device memory in use. Single-query TPU timings
therefore measure one v5e chip (Section~\ref{sec:rh-onechip}), while
cost is billed for all eight (Section~\ref{sec:cost}). The backends are
also not equally isolated: the TPU ran as a dedicated Kubernetes
Job, whereas the CPU and GPU ran on shared Colab runtimes whose
underlying machine is chosen by the provider and may change between
sessions. Section~\ref{sec:rh-variability} reports how repeated CPU and
GPU timings varied across sessions; we did not isolate a cause, and
shared tenancy remains a hypothesis there rather than a demonstrated
one.

\subsection{Software}
\label{sec:meth-software}

All AlphaFold2 runs use the public DeepMind implementation, cloned from
its default branch at run time. The August runs did not record which
commit that resolved to; the September reruns do, in their environment
files. The AlphaFold2 source is not modified. One compatibility shim
is applied from outside it: our scripts wrap \texttt{jax.numpy.clip} so
that the \texttt{a\_min}/\texttt{a\_max} keyword arguments used by
AlphaFold2 map onto the argument names of current JAX releases. No
container image digest or lockfile was used for the original runs: TPU
Jobs started from a floating \texttt{python:3.12-slim} image and
installed packages at start-up, with only JAX pinned, and the Colab
notebooks installed JAX and AlphaFold2's dependencies unpinned. For the
September reruns, each result is stored together with the benchmark
repository commit it was run from and a capture of the environment
(installed package versions and host CPU).

\subsection{Workload}
\label{sec:meth-workload}

Every AlphaFold2 run uses the \texttt{model\_3} configuration (no
templates), a single ensemble member, no recycling
(\texttt{num\_recycle}~$=0$) and float32 precision, unless an experiment
varies one of these explicitly (the recycling, precision and
model-comparison sweeps).

\paragraph{Input sequences.} Single-query experiments fold a fixed
synthetic 118-residue sequence. The sequence-length sweep, the
\texttt{vmap} batching experiment, the multi-query \texttt{pmap}
experiment and the chip-count scaling grid instead use the 20-letter
amino-acid alphabet repeated to the target length (rotated per batch
item), including at 118 residues. These two input families are
therefore not identical. The one place the paper sets a number from
each family side by side is the 6.92$\times$ multi-query throughput
gain of Section~\ref{sec:parallelism}, whose denominator is the
fixed-sequence single-query baseline; we flag that mismatch where the
figure appears, and treat the matched chip-count grid as the cleaner
comparison.

\paragraph{MSA and features.} To isolate model execution from the data
pipeline, no genetic database search is run and no templates are used:
the MSA contains only the query sequence, built with AlphaFold2's own
parsing and feature code. Feature processing happens before any timer
starts. The trivial MSA does not shrink the compiled graph, because
AlphaFold2 pads its MSA tensors to fixed sizes: the logged model inputs
are \texttt{msa\_feat} of shape $(1, 512, 118, 49)$ and
\texttt{extra\_msa} of shape $(1, 5120, 118)$. All feature, parameter
and prediction seeds are fixed to 0, with one exception: the external
ensemble experiment of Section~\ref{sec:parallelism} builds its eight
feature sets with seeds 0--7, since identical seeds would give
identical members. Parameter initialization and the shared model PRNG
stay at 0 there too.

\paragraph{Why random weights.} AlphaFold2 is run with randomly
initialized parameters; the trained checkpoint is never loaded. This is a
systems study: the question is where time goes during compilation and
execution, not how accurate the predicted structures are, and the
compiled graph's shapes and operations do not depend on parameter
values. The predicted structures are therefore meaningless and are
never interpreted, and timing equivalence with trained weights is
argued rather than measured: we did not time a run with trained
weights, so value-dependent runtime effects cannot be ruled out.
The AlphaFold3 runs (Section~\ref{sec:alphafold3}) use the real, trained
AlphaFold3 weights.

\subsection{Timing procedure}
\label{sec:meth-timing}

All times are host wall-clock (\texttt{time.time()}). Every timed
\texttt{predict} call is followed by \texttt{jax.block\_until\_ready},
so that asynchronous accelerator work is included; the
\texttt{init\_params} timer has no such explicit barrier, so any device
work it leaves outstanding is not counted in that region. Each single-query process
times three consecutive regions: parameter initialization
(\texttt{init\_params}, including its own tracing and compilation), a
first \texttt{predict} call, which compiles and executes the model, and
a second, identical \texttt{predict} call, which reuses the compiled
program and is reported as the \emph{steady state}. The first call is
the only warm-up. The timed \texttt{predict} region also includes the
confidence metrics that AlphaFold2 computes on the host after the
forward pass.

\paragraph{Repeated steady-state calls.} In the original August
baselines, the steady state is a single call per process. For the
September CPU and GPU reruns, the committed script repeats the
steady-state call a configurable number of times and makes the profiler
optional: each rerun times five consecutive steady-state calls in the
same process, with the profiler disabled, and we report their mean and
sample standard deviation.

\paragraph{Profiler overhead on the first call.} In the August and TPU
runs of the single-query and \texttt{vmap} scripts, the first
\texttt{predict} runs inside a
\texttt{jax.profiler.trace} context, and closing that context is inside
the timer. The Colab logs let us measure this directly: after
\texttt{predict} returned, the timer kept running for another 36.00\,s
on CPU and 42.02\,s on GPU, while the same \texttt{block\_until\_ready}
took under a millisecond on the second call. Reported first-call times
from these scripts therefore include trace finalization; steady-state
times are unaffected. Section~\ref{sec:rh-coldstart} reports cold-start
ratios corrected for this overhead on CPU and GPU. No TPU run log
survives, so the size of the overhead on TPU is unknown. The
\texttt{pmap}, ensemble-sharding and auto-mesh scripts do not use the
profiler, so their first-call times are not comparable with those of
the single-query script.

\paragraph{Multi-chip experiments.} The multi-query \texttt{pmap} and
\texttt{vmap} experiments report throughput as batch size divided by the
steady-state call time. Device memory is read once per process from
JAX's per-device memory statistics after the second call; the reported
peak is the process-lifetime peak.

\paragraph{AlphaFold3.} AlphaFold3 is timed from its own log: model
inference for one seed and five diffusion samples, which includes JIT
compilation because AlphaFold3 makes a single model call per process
with no warm-up. Because recycling depth, compilation and sample
accounting all differ between the two models, we report no
AlphaFold3-to-AlphaFold2 speed ratio (Section~\ref{sec:alphafold3}).

\subsection{Repetitions and noise}
\label{sec:meth-repetitions}

Most experiments are a single process per configuration point. The
auto-mesh partitioning experiment was run twice, with steady states of
0.472\,s and 0.473\,s, but the only
deliberate \emph{noise characterization} on the TPU ran the
single-query benchmark three
times in the same pod, with a coefficient of variation of 0.12\% on
steady state, 0.82\% on \texttt{init\_params} and 1.00\% on the first
call. This noise estimate applies to the TPU setup only. Access to the
TPU cluster ended with the course, so the TPU measurements cannot be
repeated. For the shared Colab runtimes, we instead reran the CPU
baseline in three further sessions (2026-09-15, -16 and -17) and the GPU
baseline in one (2026-09-15), five steady-state calls each
(Section~\ref{sec:rh-variability}). Experiments that share a Kubernetes
Job ran back to back in one pod, in a fixed order.

\subsection{Provenance}
\label{sec:meth-provenance}

The original results were committed together with the benchmark scripts
in a single commit, so version control cannot tie a result to the
script revision that produced it. The August CPU and GPU baselines were
produced by an uncommitted copy of the script embedded in the Colab
notebooks, with the same timing structure as the committed version; the
TPU baseline result matches neither surviving script version. TPU run
dates were not recorded. Version control only bounds them from above,
since all TPU results were already committed by 2026-08-11 (most of them
on 2026-08-08). The September
reruns close this gap for CPU and GPU, since each is tied to an explicit
commit and a captured environment.

The consequence is that the August baseline runs, and the headline
ratios computed from them, cannot be regenerated from the repository
as it stands. We know what was measured and from which
result file each number comes, because the catalogue records that. What
we cannot do is re-execute the exact configuration that produced them.
The September reruns are reproducible in that stronger sense, and the
comparison between the two is reported in
Section~\ref{sec:rh-variability}. Everything else our records do not
contain is listed with its consequences in Section~\ref{sec:limitations}.

